\section{最速降线}
\markboth{14. 最速降线}{14. 最速降线}

欧洲的数学界在十七世纪时盛行一种挑战的风气：一个人公开提出一个或一些数学难题，大家都来做，看谁做得快、做得对、做得好。1696年，瑞士数学家约翰·贝努利提出一个难题向全欧洲的数学家挑战，题目的大意是：设在竖直平面内有一条曲线，一个质点由于重力的作用，从这条曲线的较高的端点沿曲线滑到较低的端点，问这条曲线是什么形状时，滑行所需的时间最短？（摩擦和空气阻力都忽略）这就是著名的“最速降线”问题。

设两个端点在竖直平面内的直角坐标是 $(0, 0)$ 和 $(a, b)$，曲线方程是 $y = f(x)$。曲线要通过这两点，所以 $f(0) = 0$，$f(a) = b$。根据第 13 节中导出的公式，沿曲线 $y = f(x)$ 滑行所需的时间是
\[
  T = \int_{0}^{a} \frac{\sqrt{1 + y'^2}}{\sqrt{2gy}}\dd x.
\]
问题归结为求函数 $y = f(x)$，满足条件 $f(0) = 0$，$f(a) = b$，并且使由上面的定积分得到的数 $T$ 取最小值。

这个问题看上去好象是个很普通的极值问题，但只要仔细一分析，就知道它非同寻常。人们熟知的极值问题，是函数已经确定，再求函数的极大值或极小值。而现在是不知道函数的表达式，倒过来要求出一种函数，使得整个式子达到极小。很显然，这是一种崭新的问题。因而用当时所知道的各种数学工具，包括微分法在内，都难以对付这种陌生而复杂的极值问题。

经过一番努力，这个问题还是有几个人解出来了；解答者除去挑战人约翰·贝努利自己之外，还有牛顿、莱布尼兹、约翰的哥哥雅各·贝努利等。答案就是上一节中那倒放的摆线弧。沿摆线弧滑下，比任何曲线都快。

这个问题带来了意外的收获，那就是在这个解法的基础上，后来发展成为一门非常有用的数学新分支——变分法。在变分法教材里可以找到最速降线是摆线的证明。不过我们通过第 13 节中对摆线滑梯的计算，也可对摆线在最速降线问题中的优越地位略知一二了。
